% --- CORRECCION: Forzar interpretacion raw de la entrada ---
\UseRawInputEncoding
\documentclass[aspectratio=169, 10pt, xcolor=table]{beamer}

% --- Configuracion del Tema ---
\usetheme{Madrid}
\usecolortheme{dolphin}

% --- Paquetes necesarios ---
\usepackage[utf8]{inputenc}
\usepackage[T1]{fontenc}
\usepackage[spanish,es-noquoting]{babel}
\usepackage{amsmath, amssymb}
\usepackage{mathtools}
\usepackage{booktabs}
\usepackage{graphicx}
\usepackage{listings}
\usepackage{tikz}
\usetikzlibrary{arrows.meta, positioning, shapes.geometric, babel}
\usepackage{etoolbox}
\usepackage{upquote}

% --- Estilo para codigo SQL ---
\definecolor{codegreen}{rgb}{0,0.6,0}
\definecolor{codegray}{rgb}{0.5,0.5,0.5}
\definecolor{codepurple}{rgb}{0.58,0,0.82}
\definecolor{backcolour}{rgb}{0.95,0.95,0.92}
\lstdefinestyle{mysql}{
    language=SQL,
    backgroundcolor=\color{backcolour},
    commentstyle=\color{codegreen},
    keywordstyle=\color{blue},
    numberstyle=\tiny\color{codegray},
    stringstyle=\color{codepurple},
    basicstyle=\ttfamily\scriptsize,
    breaklines=true,
    frame=single,
    showstringspaces=false,
    literate={á}{{\'a}}1 {é}{{\'e}}1 {í}{{\'i}}1 {ó}{{\'o}}1 {ú}{{\'u}}1 {ñ}{{\~n}}1 {ü}{{\"u}}1
}
\lstdefinestyle{python}{
    language=Python,
    backgroundcolor=\color{backcolour},
    commentstyle=\color{codegreen},
    keywordstyle=\color{blue},
    numberstyle=\tiny\color{codegray},
    stringstyle=\color{codepurple},
    basicstyle=\ttfamily\scriptsize,
    breaklines=true,
    frame=single,
    showstringspaces=false,
    literate={á}{{\'a}}1 {é}{{\'e}}1 {í}{{\'i}}1 {ó}{{\'o}}1 {ú}{{\'u}}1 {ñ}{{\~n}}1 {ü}{{\"u}}1
}
\lstset{style=mysql} % Por defecto SQL, pero se puede cambiar con \lstset{style=python}

% --- Pie de pagina personalizado ---
\makeatletter
\setbeamertemplate{footline}{
  \leavevmode%
  \hbox{%
  \begin{beamercolorbox}[wd=.333333\paperwidth,ht=2.25ex,dp=1ex,center]{author in head/foot}%
    \usebeamerfont{author in head/foot}\insertshortauthor
  \end{beamercolorbox}%
  \begin{beamercolorbox}[wd=.333333\paperwidth,ht=2.25ex,dp=1ex,center]{title in head/foot}%
    \usebeamerfont{title in head/foot}Sesión 8
  \end{beamercolorbox}%
  \begin{beamercolorbox}[wd=.333333\paperwidth,ht=2.25ex,dp=1ex,right]{date in head/foot}%
    \usebeamerfont{date in head/foot}BBDD \hfill \insertframenumber/\inserttotalframenumber \hspace*{0.3cm}
  \end{beamercolorbox}}%
  \vskip0pt%
}
\makeatother

% --- Datos de la portada ---
\title[SESIÓN 8 BBDD]{\textbf{Sesión 8}}
\subtitle{Captura de Datos y Extracción Masiva}
\author[Prof. C. Sánchez]{Carlos César Sánchez Coronel}
\institute{\textbf{BASES DE DATOS}}
\date{}

\setbeamertemplate{navigation symbols}{}

\AtBeginSection[]{
  \begin{frame}
    \frametitle{Contenido de la sección}
    \tableofcontents[currentsection]
  \end{frame}
}

\begin{document}

\begin{frame}
    \titlepage
\end{frame}

\begin{frame}{Contenido general}
    \tableofcontents
\end{frame}

% ============================================================================
% SECCIÓN 1: INTRODUCCIÓN A LA CAPTURA DE DATOS
% ============================================================================
\section{Introducción a la Captura de Datos}

\begin{frame}{Importancia de la Captura de Datos}
    \begin{itemize}
        \item \textbf{Alimentación de modelos de IA}: Los modelos requieren grandes volúmenes de datos diversos (texto, imágenes, sensores).
        \item \textbf{Toma de decisiones en tiempo real}: Fraude, recomendaciones, monitoreo.
        \item \textbf{Análisis histórico}: Tendencias, estudios retrospectivos.
        \item \textbf{Integración de sistemas}: Unificar fuentes heterogéneas.
    \end{itemize}
\end{frame}

\begin{frame}{Clasificación de Fuentes de Datos}
    \begin{block}{Según origen}
        \begin{itemize}
            \item \textbf{Internas}: BD transaccionales, logs, archivos corporativos.
            \item \textbf{Externas}: Web, APIs, datos gubernamentales, redes sociales.
        \end{itemize}
    \end{block}
    \begin{block}{Según forma de entrega}
        \begin{itemize}
            \item \textbf{Batch}: Datos en bloques periódicos (diario, semanal).
            \item \textbf{Streaming}: Datos continuos en tiempo real (eventos, sensores).
        \end{itemize}
    \end{block}
\end{frame}

% ============================================================================
% SECCIÓN 2: WEB SCRAPING
% ============================================================================
\section{Fuentes Externas y Web Scraping}

\begin{frame}{Web Scraping}
    \begin{block}{Definición}
        Extracción automatizada de información de sitios web.
    \end{block}
    \begin{exampleblock}{Casos de uso}
        \begin{itemize}
            \item Entrenamiento de LLMs (Common Crawl).
            \item Monitoreo de precios de la competencia.
            \item Agregadores de noticias.
            \item Investigación académica.
        \end{itemize}
    \end{exampleblock}
\end{frame}

\begin{frame}[fragile]{Herramientas de Scraping en Python}
    \begin{description}
        \item[requests] Peticiones HTTP.
        \item[BeautifulSoup] Parseo de HTML.
        \item[Selenium] Automatización de navegadores (JavaScript).
        \item[Scrapy] Framework para scraping a gran escala.
    \end{description}
    \begin{block}{Ejemplo básico}
        \begin{lstlisting}[style=python]
import requests
from bs4 import BeautifulSoup

url = 'https://quotes.toscrape.com/'
response = requests.get(url)
soup = BeautifulSoup(response.text, 'html.parser')

for quote in soup.select('.quote'):
    text = quote.select_one('.text').get_text()
    author = quote.select_one('.author').get_text()
    print(f'{author}: {text}')
        \end{lstlisting}
    \end{block}
\end{frame}

\begin{frame}{Consideraciones Legales y Éticas}
    \begin{itemize}
        \item Respetar \texttt{robots.txt}.
        \item No sobrecargar servidores (rate limiting, delays).
        \item Verificar términos de servicio.
        \item Cumplir normativas de protección de datos (GDPR, CCPA).
    \end{itemize}
\end{frame}

\begin{frame}{Common Crawl}
    \begin{block}{¿Qué es?}
        Organización sin fines de lucro que realiza rastreos masivos de la web y pone los datos a disposición pública (petabytes en formato WARC).
    \end{block}
    \begin{alertblock}{Importancia para IA}
        Base de entrenamiento de modelos como GPT, BERT, etc. por su diversidad y escala.
    \end{alertblock}
\end{frame}

% ============================================================================
% SECCIÓN 3: APIs
% ============================================================================
\section{APIs (Application Programming Interfaces)}

\begin{frame}{Tipos de APIs}
    \begin{description}
        \item[REST] Basada en HTTP, verbos (GET, POST), devuelve JSON/XML.
        \item[GraphQL] Cliente especifica campos necesarios.
        \item[SOAP] Protocolo pesado basado en XML, entornos empresariales.
        \item[WebSockets] Comunicación bidireccional en tiempo real.
    \end{description}
\end{frame}

\begin{frame}[fragile]{Autenticación y Rate Limiting}
    \begin{itemize}
        \item \textbf{API Key} en header o parámetro.
        \item \textbf{OAuth} tokens.
        \item \textbf{Rate limiting}: límite de peticiones por tiempo; respetar para no ser bloqueado.
    \end{itemize}
    \begin{block}{Ejemplo con API REST (JSONPlaceholder)}
        \begin{lstlisting}[style=python]
import requests
import json

url = 'https://jsonplaceholder.typicode.com/posts'
response = requests.get(url)
if response.status_code == 200:
    posts = response.json()
    with open('posts.json', 'w') as f:
        json.dump(posts, f, indent=2)
else:
    print('Error:', response.status_code)
        \end{lstlisting}
    \end{block}
\end{frame}

% ============================================================================
% SECCIÓN 4: STREAMING Y CHANGE DATA CAPTURE
% ============================================================================
\section{Captura en Tiempo Real y Streaming}

\begin{frame}{Concepto de Streaming}
    \begin{block}{Procesamiento continuo de datos a medida que llegan}
        \begin{itemize}
            \item Baja latencia (ms).
            \item Eventos, sensores, logs, redes sociales.
        \end{itemize}
    \end{block}
    \begin{exampleblock}{Casos de uso}
        \begin{itemize}
            \item Detección de fraude en transacciones.
            \item Monitoreo IoT.
            \item Análisis de tendencias en redes sociales.
        \end{itemize}
    \end{exampleblock}
\end{frame}

\begin{frame}{Change Data Capture (CDC)}
    \begin{block}{Captura de cambios en bases de datos en tiempo real}
        \begin{itemize}
            \item Inserciones, actualizaciones, eliminaciones.
            \item Sin afectar el rendimiento transaccional.
        \end{itemize}
    \end{block}
    \begin{description}
        \item[Basado en logs] Lee log transaccional (WAL, binlog). Ej: Debezium.
        \item[Basado en triggers] Impacta rendimiento.
        \item[Consultas periódicas] No es tiempo real.
    \end{description}
\end{frame}

\begin{frame}{Debezium y Kafka}
    \begin{center}
        \begin{tikzpicture}[node distance=2cm, auto]
            \node (db) [draw, rectangle] {Base de Datos};
            \node (debezium) [draw, rectangle, right of=db, xshift=1.5cm] {Debezium};
            \node (kafka) [draw, rectangle, right of=debezium, xshift=1.5cm] {Kafka};
            \draw[->] (db) -- (debezium) node[midway, above] {log};
            \draw[->] (debezium) -- (kafka) node[midway, above] {mensajes};
        \end{tikzpicture}
    \end{center}
    \begin{itemize}
        \item Debezium se conecta a la BD y publica cambios en tópicos Kafka.
        \item Consumidores (Spark, aplicaciones) procesan los eventos.
    \end{itemize}
\end{frame}

\begin{frame}{Apache Kafka}
    \begin{block}{Plataforma de streaming distribuida}
        \begin{itemize}
            \item \textbf{Tópicos}: categorías de mensajes.
            \item \textbf{Productores}: publican mensajes.
            \item \textbf{Consumidores}: leen mensajes.
            \item \textbf{Partitions}: paralelismo y orden.
        \end{itemize}
    \end{block}
    \begin{alertblock}{Ventajas}
        Alta escalabilidad, persistencia, tolerancia a fallos, procesamiento en tiempo real.
    \end{alertblock}
\end{frame}

\begin{frame}{Otras Plataformas de Streaming}
    \begin{itemize}
        \item \textbf{AWS Kinesis}: servicio gestionado en AWS.
        \item \textbf{Google Pub/Sub}: global, push/pull.
        \item \textbf{Apache Pulsar}: almacenamiento en capas.
        \item \textbf{RabbitMQ}: orientado a colas (menos escalable).
    \end{itemize}
\end{frame}

% ============================================================================
% SECCIÓN 5: FORMATOS DE DATOS Y ALMACENAMIENTO
% ============================================================================
\section{Formatos de Datos y Almacenamiento}

\begin{frame}{Formatos de Archivo}
    \begin{table}
        \centering
        \begin{tabular}{l|l|l|l}
            \toprule
            Formato & Tipo & Compresión & Uso típico \\
            \midrule
            CSV & Texto fila & Baja & Intercambio simple \\
            JSON & Texto semi-estructurado & Media & APIs, documentos \\
            Avro & Binario fila con esquema & Alta & Streaming (Kafka) \\
            Parquet & Columnar & Alta & Análisis, data lakes \\
            ORC & Columnar & Alta & Hive, Hadoop \\
            WARC & Archivo web & Variable & Common Crawl \\
            \bottomrule
        \end{tabular}
    \end{table}
\end{frame}

\begin{frame}{Data Lakes vs Data Warehouses}
    \begin{block}{Data Lake}
        \begin{itemize}
            \item Almacena datos en bruto (formatos nativos).
            \item Flexible, estructurados y no estructurados.
            \item Ej: Amazon S3, Azure Data Lake, HDFS.
        \end{itemize}
    \end{block}
    \begin{block}{Data Warehouse}
        \begin{itemize}
            \item Datos procesados y estructurados (esquema estrella).
            \item Optimizado para consultas analíticas rápidas.
            \item Ej: Snowflake, BigQuery, Redshift.
        \end{itemize}
    \end{block}
\end{frame}

\begin{frame}{Particionamiento en Data Lakes}
    \begin{block}{Organización en carpetas}
        \texttt{s3://bucket/ventas/year=2025/month=03/day=15/ventas.parquet}
    \end{block}
    \begin{itemize}
        \item Permite a motores como Spark leer solo las particiones necesarias.
        \item Reduce E/S y acelera consultas.
    \end{itemize}
\end{frame}

% ============================================================================
% SECCIÓN 6: PROCESAMIENTO A GRAN ESCALA
% ============================================================================
\section{Procesamiento a Gran Escala}

\begin{frame}{Apache Spark}
    \begin{block}{Motor unificado de procesamiento distribuido}
        \begin{itemize}
            \item \textbf{Spark SQL}: consultas SQL sobre DataFrames.
            \item \textbf{DataFrames}: API similar a tablas.
            \item \textbf{Structured Streaming}: procesamiento en streaming.
            \item \textbf{MLlib}: machine learning a escala.
        \end{itemize}
    \end{block}
\end{frame}

\begin{frame}[fragile]{Ejemplo con PySpark}
    \begin{block}{Leer CSV, calcular total por producto y guardar en Parquet}
        \begin{lstlisting}[style=python]
from pyspark.sql import SparkSession

spark = SparkSession.builder.appName("ventas").getOrCreate()

df = spark.read.option("header", "true").csv("ruta/ventas.csv")
df = df.withColumn("monto", df["cantidad"].cast("int") * df["precio"].cast("float"))
resultado = df.groupBy("producto").sum("monto").withColumnRenamed("sum(monto)", "total")
resultado.write.parquet("resultado.parquet")
        \end{lstlisting}
    \end{block}
\end{frame}

\begin{frame}{Apache Flink vs Spark Streaming}
    \begin{description}
        \item[Spark Streaming] Micro-batches (intervalos de tiempo). Robusto, maduro.
        \item[Flink] Streaming puro, evento por evento, menor latencia. Ideal para fraudes, monitoreo.
    \end{description}
\end{frame}

\begin{frame}{Hadoop MapReduce}
    \begin{itemize}
        \item Modelo de programación distribuida (map y reduce).
        \item Lento porque escribe resultados intermedios en disco.
        \item Reemplazado por Spark en la mayoría de casos.
    \end{itemize}
\end{frame}

% ============================================================================
% SECCIÓN 7: EJEMPLO INTEGRADOR: PIPELINE DE TWEETS
% ============================================================================
\section{Ejemplo Integrador: Captura de Tweets en Tiempo Real}

\begin{frame}{Arquitectura del Pipeline}
    \begin{enumerate}
        \item \textbf{Fuente}: API de Twitter (streaming) con autenticación OAuth.
        \item \textbf{Ingesta}: Productor Kafka (Python) publica tweets en tópico \texttt{tweets}.
        \item \textbf{Procesamiento}: Spark Structured Streaming lee de Kafka, limpia (quita RTs), extrae hashtags, y guarda en Parquet particionado por fecha.
        \item \textbf{Almacenamiento}: Data lake en S3 (Parquet).
        \item \textbf{Análisis batch}: Spark genera reportes diarios.
    \end{enumerate}
\end{frame}

\begin{frame}[fragile]{Productor Kafka en Python}
    \begin{lstlisting}[style=python]
from kafka import KafkaProducer
import tweepy, json

producer = KafkaProducer(
    bootstrap_servers='localhost:9092',
    value_serializer=lambda v: json.dumps(v).encode('utf-8')
)

class StreamListener(tweepy.Stream):
    def on_data(self, raw_data):
        tweet = json.loads(raw_data)
        producer.send('tweets', value=tweet)
        return True

stream = StreamListener(consumer_key, consumer_secret, access_token, access_token_secret)
stream.filter(track=['#datascience'])
    \end{lstlisting}
\end{frame}

\begin{frame}[fragile]{Consumidor Spark Structured Streaming}
    \begin{lstlisting}[style=python]
from pyspark.sql import SparkSession
from pyspark.sql.functions import from_json, col
from pyspark.sql.types import StructType, StructField, StringType, TimestampType

spark = SparkSession.builder.appName("tweet_analysis").getOrCreate()

schema = StructType([
    StructField("id", StringType()),
    StructField("text", StringType()),
    StructField("created_at", TimestampType())
])

df = spark.readStream.format("kafka") \
    .option("kafka.bootstrap.servers", "localhost:9092") \
    .option("subscribe", "tweets") \
    .load() \
    .select(from_json(col("value").cast("string"), schema).alias("data")) \
    .select("data.*")

# Limpieza: quitar retweets
df_clean = df.filter(~col("text").startswith("RT"))

# Escribir en Parquet particionado por fecha
query = df_clean.writeStream \
    .format("parquet") \
    .partitionBy("created_at") \
    .option("path", "s3a://bucket/tweets/") \
    .option("checkpointLocation", "/tmp/checkpoint") \
    .trigger(processingTime="5 minutes") \
    .start()

query.awaitTermination()
    \end{lstlisting}
\end{frame}

% ============================================================================
% SECCIÓN 8: EJERCICIOS RESUELTOS (DEL SOLUCIONARIO)
% ============================================================================
\section{Ejercicios Resueltos}

\begin{frame}[fragile]{Ejercicio 1: Scraping de una Página Web}
    \textbf{Enunciado:} Extraer citas, autores y etiquetas de \url{http://quotes.toscrape.com} y guardar en CSV.
    \begin{block}{Solución}
        \begin{lstlisting}[style=python]
import requests, csv
from bs4 import BeautifulSoup

url = 'http://quotes.toscrape.com'
quotes = []
while url:
    response = requests.get(url)
    soup = BeautifulSoup(response.text, 'html.parser')
    for quote in soup.select('.quote'):
        text = quote.select_one('.text').get_text()
        author = quote.select_one('.author').get_text()
        tags = [tag.get_text() for tag in quote.select('.tag')]
        quotes.append([text, author, ', '.join(tags)])
    next_btn = soup.select_one('.next a')
    url = next_btn['href'] if next_btn else None

with open('quotes.csv', 'w', newline='', encoding='utf-8') as f:
    writer = csv.writer(f)
    writer.writerow(['quote', 'author', 'tags'])
    writer.writerows(quotes)
        \end{lstlisting}
    \end{block}
\end{frame}

\begin{frame}[fragile]{Ejercicio 2: Consumo de API Pública}
    \textbf{Enunciado:} Obtener posts de JSONPlaceholder y guardarlos en JSON.
    \begin{block}{Solución}
        \begin{lstlisting}[style=python]
import requests, json

response = requests.get('https://jsonplaceholder.typicode.com/posts')
posts = response.json()
with open('posts.json', 'w') as f:
    json.dump(posts, f, indent=2)
        \end{lstlisting}
    \end{block}
\end{frame}

\begin{frame}[fragile]{Ejercicio 3: Simulación de CDC con Triggers en PostgreSQL}
    \textbf{Enunciado:} Crear trigger que registre cambios en \texttt{clientes} y un script Python que lea la auditoría y publique en Kafka.
    \begin{block}{Trigger (PostgreSQL)}
        \begin{lstlisting}[style=mysql]
CREATE TABLE clientes_audit (...);
CREATE OR REPLACE FUNCTION audit_clientes() RETURNS TRIGGER AS $$
BEGIN
    IF TG_OP = 'INSERT' THEN
        INSERT INTO clientes_audit (cliente_id, nombre_nuevo, email_nuevo, operacion)
        VALUES (NEW.id, NEW.nombre, NEW.email, 'I');
    ELSIF TG_OP = 'UPDATE' THEN
        INSERT INTO clientes_audit (cliente_id, nombre_anterior, nombre_nuevo, email_anterior, email_nuevo, operacion)
        VALUES (OLD.id, OLD.nombre, NEW.nombre, OLD.email, NEW.email, 'U');
    ELSIF TG_OP = 'DELETE' THEN
        INSERT INTO clientes_audit (cliente_id, nombre_anterior, email_anterior, operacion)
        VALUES (OLD.id, OLD.nombre, OLD.email, 'D');
    END IF;
    RETURN NULL;
END;
$$ LANGUAGE plpgsql;
        \end{lstlisting}
    \end{block}
\end{frame}

\begin{frame}[fragile]{Script Python para leer auditoría y enviar a Kafka}
    \begin{lstlisting}[style=python]
import psycopg2
from kafka import KafkaProducer
import json

conn = psycopg2.connect("dbname=test user=postgres")
cur = conn.cursor()
cur.execute("SELECT * FROM clientes_audit WHERE fecha > NOW() - INTERVAL '1 minute'")
rows = cur.fetchall()

producer = KafkaProducer(bootstrap_servers='localhost:9092',
                         value_serializer=lambda v: json.dumps(v).encode('utf-8'))

for row in rows:
    data = {'id': row[0], 'cliente_id': row[1], 'operacion': row[7], 'fecha': str(row[6])}
    producer.send('clientes_audit', value=data)
    \end{lstlisting}
\end{frame}

\begin{frame}[fragile]{Ejercicio 4: Procesamiento Batch con Spark}
    \textbf{Enunciado:} Cargar CSV de ventas, calcular total por producto y guardar en Parquet.
    \begin{block}{Solución}
        \begin{lstlisting}[style=python]
from pyspark.sql import SparkSession

spark = SparkSession.builder.appName("ventas").getOrCreate()
df = spark.read.option("header", "true").csv("ruta/ventas*.csv")
df = df.withColumn("monto", df["cantidad"].cast("int") * df["precio"].cast("float"))
resultado = df.groupBy("producto").sum("monto").withColumnRenamed("sum(monto)", "total")
resultado.write.parquet("resultado.parquet")
        \end{lstlisting}
    \end{block}
\end{frame}

% ============================================================================
% SECCIÓN 9: GLOSARIO
% ============================================================================
\section{Glosario}

\begin{frame}{Términos Clave}
    \begin{description}
        \item[Web Scraping] Extracción automatizada de datos web.
        \item[API] Interfaz para comunicación entre sistemas.
        \item[CDC] Change Data Capture, captura de cambios en BD.
        \item[Streaming] Procesamiento continuo de datos.
        \item[Kafka] Plataforma de streaming distribuida.
        \item[Debezium] Herramienta de CDC sobre Kafka.
        \item[Parquet] Formato columnar optimizado para análisis.
        \item[Data Lake] Almacén de datos en bruto.
        \item[Spark] Motor de procesamiento distribuido.
    \end{description}
\end{frame}

% --- Diapositiva final ---
\begin{frame}
    \centering

    
    \vspace{1cm}
    \large ¡Gracias!
\end{frame}

\end{document}